波长 \(\lambda\) 的准单色空间非相干扩展恒星，在天空上可近似为均匀亮度圆盘，角直径 \(\delta\ll1\) 弧度。两个小孔干涉望远镜基线长度 \(B\) 沿天空平面方向，且恒星角大小足够小可取平面角坐标；采用范西特--泽尼克定理，忽略大气湍流和单孔衍射包络对整个星盘的变化。
(1) 写出归一化复空间相干度对应的二维亮度 Fourier 积分。
(2) 求均匀圆盘的相干度和可见度解析式。
(3) 求第一个可见度零点所对应的基线与恒星角直径关系。
(4) 推导小基线近似，并讨论为何仅靠可见度模长不能测定源的方位角。
